\documentclass[12pt]{amsart}
%   \overfullrule=5pt
\usepackage[active]{srcltx}
\usepackage{calc,amssymb,amsthm,amsmath,amscd, eucal,ulem}
\usepackage{alltt}
%\usepackage[left=1in,top=1in,right=1in,bottom=1in]{geometry}
\synctex=1
%\usepackage{mathtools}
\RequirePackage[dvipsnames,usenames]{color}
\let\mffont=\sf
\normalem
\input{mabliautoref.sty}
\input{xy}
\xyoption{all}
\input{kmacros3.sty}
\usepackage{tikz}
\usepackage{amsfonts, mathrsfs}
\usepackage{calligra}
\usepackage{stmaryrd} %power series bracket
%\usepackage{dcpic, pictexwd}
%\usepackage{pdfsync}
\usepackage[left=1.1in,top=1in,right=1.1in,bottom=1in]{geometry}
\usepackage{enumitem}

\numberwithin{equation}{theorem}

%\renewcommand{\thefootnote}{\fnsymbol{footnote}}
\newcommand{\D}{\displaystyle}
\newcommand{\til}{\widetilde}
\newcommand{\ol}{\overline}
\newcommand{\F}{\mathbb{F}}
\DeclareMathOperator{\E}{E}
\renewcommand{\:}{\colon}
\newcommand{\eg}{{\itshape e.g.} }
\renewcommand{\m}{\mathfrak{m}}
\renewcommand{\n}{\mathfrak{n}}
\newcommand{\Tt}{{\mathfrak{T}}}
\newcommand{\calC}{\mathcal{C}}
\newcommand{\RamiT}{\mathcal{R_{T}}}
\DeclareMathOperator{\edim}{edim}
\DeclareMathOperator{\length}{length}
\DeclareMathOperator{\monomials}{monomials}
\DeclareMathOperator{\Fitt}{Fitt}
\DeclareMathOperator{\depth}{depth}
\newcommand{\Frob}[2]{{#1}^{1/p^{#2}}}
\newcommand{\Frobp}[2]{{(#1)}^{1/p^{#2}}}
\newcommand{\FrobP}[2]{{\left(#1\right)}^{1/p^{#2}}}
\newcommand{\extends}{extends over $\Tt${}}
\newcommand{\extension}{extension over $\Tt${}}
\newcommand{\fg}{\pi_1^{\textnormal{\'{e}t}}} %Fundamental group
\DeclareMathOperator{\ns}{ns}
%\newcommand{\Spec}{\text{Spec}} %Spectrum
\newcommand{\pSpec}{\text{Spec}^{\circ}} %Puntured Spectrum
\newcommand{\etInCdOne}{\etale{} in codimension $1$}
\DeclareMathOperator{\DIV}{Div}
%\renewcommand{\char}{\charact}


\usepackage{setspace}
\usepackage{hyperref}
%\usepackage{enumerate}
\usepackage{graphicx}
\usepackage[all,cmtip]{xy}

\usepackage{verbatim}

\theoremstyle{theorem}
\newtheorem*{mainthm}{Main Theorem}
\newtheorem*{mainthma}{Main Theorem A}
\newtheorem*{mainthmb}{Main Theorem B}
\newtheorem*{mainthmc}{Main Theorem C}
\newtheorem*{atheorem}{Theorem}


\newcommand{\sol}[1]{ \vskip 6pt{\color{blue} {\bf Solution:} #1} \vskip 6pt}
%The todo box!
\def\todo#1{\textcolor{Mahogany}%
{\footnotesize\newline{\color{Mahogany}\fbox{\parbox{\textwidth-15pt}{\textbf{todo:
} #1}}}\newline}}

%What is going on?  Why isn't \O a script O?!?
\renewcommand{\O}{\mathscr O}

\begin{document}
\title{Exercises for Characteristic $p$ commutative algebra\\February 3rd, 2017}
\author{Due, February 13th, 2017}
%\address{Department of Mathematics\\ University of Utah\\ Salt Lake City\\ UT 84112}
%\email{schwede@math.utah.edu}


%  \subjclass[2010]{14F18, 13A35, 14B05}

\maketitle

%\section{Reduction to characteristic $p > 0$}
%\chapter{Frobenius and Kunz's theorem}
%\bibliographystyle{skalpha}
%\bibliography{MainBib}

\begin{enumerate}[label=(\arabic*)]
\item If we have a morphism of complexes, $A^{\mydot} \xrightarrow{\alpha} B^{\mydot}$, then we can always take the cone $C(\alpha)^{\mydot} = C^{\mydot} =  A[1]^{\mydot} \oplus B^{\mydot}$ with differential
\[
C^i = A^{i+1} \oplus B^i \xrightarrow{-d_A^{i+1}, \alpha^i + d_B^i} A^{i+2} \oplus B^{i+1}.
\]
Verify that this really is a complex.
\item Suppose that $0 \to A^{\mydot} \xrightarrow{\alpha} B^{\mydot} \xrightarrow{\beta} D^{\mydot} \to 0$ is a short exact sequence of complexes.  Show that $D^{\mydot}$ is quasi-isomorphic to the cone $C(\alpha)^{\mydot}$.  
\item Consider the ring $R = {\overline \bF}_p[x]$.  Show that the perfection $R^{\infty}$ of $R$ is ${\overline \bF}_p[x, x^{1/p}, x^{1/p^2}, \ldots]$ and that the map $\Spec R^{\infty} \to \Spec R$ is a bijection.
\item More generally, suppose that $R$ is a Noetherian ring and $R^{\infty}$ is its perfection (which as we've seen is usually not Noetherian).  Is $\Spec R^{\infty} \to \Spec R$ always a bijection?
\item Suppose $R$ is an $F$-finite Noetherian ring.  Show that $R$ is $F$-split if and only if $R_{\fram}$ is $F$-split for every maximal ideal $\fram \subseteq R$.
\item Suppose we have an extension of rings $R \subseteq S$ and that $M$ is an $S$-module.  Suppose that $\Hom_R(S, R) \cong S$ as $S$-modules.  We have a map
    \[
        \Hom_S(M, S) \times \Hom_R(S, R) \to \Hom_R(M, R)
    \]
    by composition $(\alpha, \beta) \mapsto \beta \circ \alpha$.  Show that this map is surjective.
\item Suppose that we have inclusions $R \subseteq S \subseteq T$ of rings and that $\Hom_R(S, R) \cong_S S$ and that $\Hom_S(T,S) \cong_T T$.  Choose $\Phi \in \Hom_R(S, R)$ which generates the $\Hom$-module as an $S$-module and $\Psi \in \Hom_S(T, S)$ generating the $\Hom$-module as a $T$-module.  Prove that $\Phi \circ \Psi$ generates $\Hom_R(T, R)$ as a $T$-module.
\end{enumerate}



\end{document}


